% !TeX program = xelatex
\documentclass[11pt,a4paper]{article}
\usepackage[margin=23mm]{geometry}
\usepackage{fontspec}
\setmainfont{Latin Modern Roman}
\usepackage{amsmath,amssymb,mathtools}
\usepackage{xcolor,enumitem,booktabs,longtable,array,xurl,hyperref}
\definecolor{muted}{HTML}{53636C}
\hypersetup{colorlinks=true,urlcolor=blue,linkcolor=blue}
\setlength{\parindent}{0pt}
\setlength{\parskip}{5pt}
\newcommand{\R}{\mathbb R}
\newcommand{\E}{\mathbb E}
\newcommand{\N}{\mathbb N}
\newcommand{\ind}{\mathbf 1}
\newcommand{\Var}{\operatorname{Var}}
\newcommand{\Cov}{\operatorname{Cov}}
\newcommand{\supp}{\operatorname{supp}}
\DeclareMathOperator*{\argmin}{arg\,min}
\DeclareMathOperator*{\argmax}{arg\,max}
\newcommand{\entrymeta}[3]{{\sffamily\small\color{muted}\textbf{\texttt{#1}}\quad|\quad #2\par #3\par}}
\newcommand{\Problemref}[1]{\texttt{#1}}
\newcommand{\JELref}[2]{\texttt{#2} (JEL \texttt{#1})}
\begin{document}
\section*{Equilibrium selection in games}
\textbf{Problem ID:} \texttt{C72-489}\quad \textbf{JEL:} \texttt{C72}\par
% BEGIN ECONOMICS STATEMENT
\label{problem:OP-0001}
\label{atlas:prob:M1}

\noindent\textbf{Original atlas ID:} M1. \textbf{Formulation:} editorial specification of a research family.

\subsection*{Definitions and assumptions}

Fix a class $\mathcal G$ of finite normal-form games with players $N=\{1,\ldots,n\}$, where $n$ is a positive integer, nonempty finite action sets $A_i$, and utilities $u_i:A\to\mathbb R$, where $A=\prod_iA_i$. Here $\Delta(A_i)$ is the probability simplex on $A_i$, equipped with its Borel sigma-algebra. For game $G$, let $\mathcal N(G)\subseteq\prod_i\Delta(A_i)$ be its mixed-strategy Nash-equilibrium set: $\sigma\in\mathcal N(G)$ if every $\sigma_i$ maximizes expected $u_i$ against $\sigma_{-i}$. A context $c$ specifies the protocol, information, communication opportunities, and participant population. A selection rule is a probability kernel $q(\mathrm d\sigma\mid G,c)$ supported on $\mathcal N(G)$. This is a maintained behavioral hypothesis, not a consequence of equilibrium existence. Let $H_{G,c}(\mathrm dz\mid\sigma)$ be a specified observation kernel linking a selected equilibrium to measured session outcomes $z$. Each instance must declare measurable game and context domains and an outcome space $(\mathcal Z,\mathscr Z)$, with $B\in\mathscr Z$. If player or action sets vary across games, the instance must also supply a common measurable realization of the game-indexed strategy spaces on which the kernels are defined. The equalities below are required at every observed pair $(G,c)$, not merely almost everywhere under an unstated design law.

\subsection*{Formal research question}

Given population outcome laws $P_{G,c}$ and a prespecified restricted family $\mathcal Q$ of selection kernels, characterize
\[
\begin{aligned}
\mathcal Q(P)=\bigl\{q\in\mathcal Q:\;&
 P_{G,c}(B)=\int H_{G,c}(B\mid\sigma)q(\mathrm d\sigma\mid G,c)\\
 &\text{for all observed }(G,c)\text{ and measurable }B\bigr\}.
\end{aligned}
\]
Under what experimentally enforceable restrictions is this set a singleton, and which predictions for unobserved games or contexts are common to all its elements? For a specified unobserved pair $(G,c)$, define the compatible observable-law set
\[
\mathcal P_{G,c}(P)=\left\{\int H_{G,c}(\cdot\mid\sigma)q(\mathrm d\sigma\mid G,c):q\in\mathcal Q(P)\right\}.
\]
If $\mathcal Q(P)$ is empty, the maintained model is incompatible with the population laws and supplies no compatible prediction; agreement over an empty set is not an identified prediction. If it is nonempty, an observable law is identified when $\mathcal P_{G,c}(P)$ is a singleton; a declared functional of that law is identified when it is constant on this set. Singleton and multiple-element selection sets must be distinguished from singleton and multiple-element prediction sets. Emptiness of $\mathcal Q(P)$ rejects the joint restrictions on $\mathcal Q$ and $H$, not necessarily equilibrium support. Rejection of equilibrium support under the maintained $H$ requires incompatibility even when all admissible equilibrium-supported selection kernels are allowed. Population incompatibility is distinct from a finite-sample statistical rejection. When observations reject equilibrium support, enlarge the model explicitly to approximate equilibria or transient behavior and test that enlargement. For a declared distribution over held-out games, estimate predictive risk with uncertainty that includes selection-rule estimation. Each empirical instance must additionally declare the prediction object, loss, context distribution or context-assignment rule, sampling experiment, sample sizes, dependence structure, time unit and observation horizon, and uncertainty criterion. It must specify which training-sample and test-outcome randomness the risk averages over. Approximate-equilibrium or transient enlargements require their own declared approximation notion, payoff units and tolerance, or dynamic observation model; none is selected here.

\subsection*{Interpretation and scope}

Identification of latent equilibrium selection is separate from prediction of observed action frequencies. Two selection kernels can be observationally equivalent when the observation kernel does not distinguish their induced outcomes. Requiring each admissible selector to choose a single equilibrium with probability one does not by itself identify which selector generated the observations or its held-out predictions; conversely, an observable prediction can be identified without identifying the selector. Context restrictions must say whether participant composition, payoff scaling, and information precision may change. It remains to be declared which information and communication changes preserve the stated normal-form game, and which require a different strategic representation and a correspondingly specified equilibrium-support condition. The present definition of $\mathcal N(G)$ alone does not specify such an extension. Cross-game extrapolation requires shared parameters or invariance assumptions that can be tested on strategically related environments. Each instance must declare its observed design, admissible context interventions and cross-environment restrictions; randomization of contexts does not by itself enforce behavioral invariance. A useful study randomizes these contextual variables and evaluates the resulting rule on games excluded from estimation; it also records finite-horizon departures from stationarity.

\subsection*{Source audit and status}

Nash (1951) establishes existence in finite games, not empirical selection. Harsanyi--Selten (1988) supplies a theoretical selection construction. Carlsson--van Damme (1993) gives selection results for particular noisy-information perturbations; the conclusions cannot be transferred to arbitrary perturbations. All three author--year references are verified through the primary records below. The atlas question is a broad predictive agenda: many specified selection problems are solved, and these old sources do not certify that one precisely stated conjecture remains open in 2026. The kernel model and identification target above are editorial formalizations; no claim is made that they were stated in these sources.

\subsection*{References}

\begin{enumerate}

\item[S1] John Nash (1951). \emph{Non-Cooperative Games}. Annals of Mathematics 54(2), 286--295. \url{https://www.jstor.org/stable/1969529}. \textit{Locator:} equilibrium definition and existence argument. \textit{Establishes:} Existence of mixed-strategy equilibrium in finite games. \textit{Does not establish:} a universal empirical equilibrium-selection rule.

\item[S2] John C. Harsanyi and Reinhard Selten (1988). \emph{A General Theory of Equilibrium Selection in Games}. MIT Press. \url{https://mitpress.mit.edu/9780262582384/a-general-theory-of-equilibrium-selection-in-games/}. \textit{Locator:} publisher description; selection construction. \textit{Establishes:} A systematic theoretical selection proposal. \textit{Does not establish:} that its selection is invariably observed.

\item[S3] Hans Carlsson and Eric van Damme (1993). \emph{Global Games and Equilibrium Selection}. Econometrica 61(5), 989--1018. \url{https://research.tilburguniversity.edu/en/publications/global-games-and-equilibrium-selection/}. \textit{Locator:} abstract and two-by-two coordination analysis. \textit{Establishes:} Selection under specified noisy private-information perturbations. \textit{Does not establish:} selection for every game or every information perturbation.

\end{enumerate}

\subsection*{Common conventions: Source mathematical conventions}

Unless an entry states otherwise, agent and firm sets are finite. State and action spaces are measurable, strategies and outcome maps are measurable, probability laws are specified, and expectations exist for the targets under discussion. Infinite-horizon discount factors lie in $(0,1)$ and discounted objectives are finite. Norms, tolerances and complexity input sizes must be specified for computational comparisons. Constants or primitives that appear locally are inputs, not empirical facts asserted by this catalogue.

\textbf{Identification.} A statistical model $m\in\mathcal M$ implies an observable law $P_m$ and target $\theta(m)$. At law $P$, its identified set is
\[
\Theta(P;\mathcal M)=\{\theta(m):m\in\mathcal M,\ P_m=P\}.
\]
Point identification at $P$ holds when this set is a singleton. Provided at least one compatible model exists, this is equivalent to $\theta(m_1)=\theta(m_2)$ for all $m_1,m_2\in\mathcal M$ with $P_{m_1}=P_{m_2}=P$. Identification throughout the model requires $P_{m_1}=P_{m_2}$ to imply $\theta(m_1)=\theta(m_2)$ for every pair $m_1,m_2\in\mathcal M$. Sharp bounds mean bounds attained, or approached when the boundary is not attained, by compatible models. Writing this set is a definition, not a solution: the substantive task is to establish informative restrictions and derive or estimate the set. An unrestricted-confounding model may give only trivial bounds.

\textbf{Inference.} Identification, estimability and valid uncertainty quantification are separate questions. Fix $\alpha\in(0,1)$. A confidence procedure $C_n$ over a declared law class $\mathcal P$ has asymptotic uniform coverage at level $1-\alpha$ when
\[
\liminf_{n\to\infty}\inf_{P\in\mathcal P}\Pr_P\{\theta(P)\in C_n\}\geq1-\alpha.
\]
For set-identified targets the coverage event must be defined separately, for example coverage of the full identified set. If an entry uses identification-indexed models instead of uniquely indexed laws, the same coverage convention is applied over those models. Pointwise limits do not establish uniform validity, and asymptotic validity does not guarantee finite-sample precision. Sample sizes, dependence structures and weak-identification sequences are part of each application.

\textbf{Counterfactuals.} $Y_i(a)$ is a potential outcome under a specified intervention, including its timing, implementation and versions. It is not $Y_i$ conditional on voluntarily choosing $a$. A full intervention vector permits interference; shorthand scalar treatment notation does not silently impose no interference. Assignment, selection, attrition, anticipation and measurement must be specified. A claim about an instrument's local effect is not automatically a population policy effect.

\textbf{Economic equilibrium.} A counterfactual model includes best responses, constraints, feasibility, market clearing, state transitions and expectations consistent with the stated information. When several equilibria exist, either a declared selector is an input or the target is set-valued. Comparisons across policy regimes must use the stated selection convention. Reduced-form relationships are not presumed invariant after an equilibrium-changing intervention.

\textbf{Optimization and welfare.} Optimization is over the stated feasible set. If attainment is not established, $\sup$ or $\inf$ and $\varepsilon$-optimal choices replace an asserted maximizer or minimizer. For example, with value $V=\sup_{a\in\mathcal A}W(a)<\infty$, an $\varepsilon$-optimal set is $\{a\in\mathcal A:W(a)\geq V-\varepsilon\}$ for $\varepsilon>0$. Benefits, costs, risk and health or environmental outcomes can be aggregated only using declared units and conversion weights. Interpersonal utility weights, the treatment of fiscal transfers and foreign profits, discounting, and the behavioral welfare criterion are explicit inputs. A comparative-static derivative additionally requires existence and appropriate differentiability of the selected counterfactual.

\textbf{Sources.} Local labels [S1], [S2], and so forth restart in every question. Locators identify the relevant abstract, model, section or result at the level actually checked; a bibliographic or abstract-level verification is not represented as inspection of an entire proof. Working papers are distinguished from later publications and changing versions. Source summaries are paraphrased. All URL checks and editorial formulations refer to the audit date stated above.
% END ECONOMICS STATEMENT
\end{document}
